% =====================================================================
%  ENTREGABLE I - Descripción del problema
%  Plantilla: IEEEtran journal (bare_jrnl_new_sample4.tex)
%  Compilador: pdfLaTeX  (IEEEtran.cls) + BibTeX (references.bib)
%
%  Las cifras y figuras del análisis exploratorio se llenan solas con los
%  archivos que genera el notebook EDA_Entregable1.ipynb: resultados_eda.tex
%  y las cuatro figuras en PDF (en la raíz del proyecto o en figs/). Si falta
%  alguno de esos archivos, su lugar aparece en ROJO.
% =====================================================================
\documentclass[letterpaper,journal]{IEEEtran}
\usepackage{amsmath,amsfonts}
\usepackage{algorithmic}
\usepackage{algorithm}
\usepackage{array}
\usepackage[caption=false,font=normalsize,labelfont=sf,textfont=sf]{subfig}
\usepackage{textcomp}
\usepackage{stfloats}
\usepackage{url}
\usepackage{verbatim}
\usepackage{graphicx}
\usepackage{booktabs}
\usepackage{cite}
\usepackage{xcolor} % solo para las marcas en rojo; puede quitarse al final
\usepackage[T1]{fontenc}
% Español: Resumen, Palabras clave, TABLA, Referencias y división silábica
% correcta. La opción es-minimal evita que babel cambie el formato de IEEEtran.
\usepackage[spanish,es-minimal,es-noshorthands,es-tabla]{babel}
\addto\captionsspanish{\renewcommand{\figurename}{Fig.}\renewcommand{\tablename}{TABLA}}
\renewcommand{\IEEEkeywordsname}{Palabras clave}
\hyphenation{op-tical net-works semi-conduc-tor IEEE-Xplore}

% Columna de ancho fijo alineada a la izquierda (tablas)
\newcolumntype{L}[1]{>{\raggedright\arraybackslash}p{#1}}
% Marcas de contenido pendiente
\newcommand{\pendiente}[1]{\textcolor{red}{[#1]}}
\newcommand{\figpendiente}[1]{\fbox{\parbox[c][2.3cm][c]{0.92\columnwidth}{\centering\footnotesize\textcolor{red}{Figura pendiente: #1}}}}
% Cifras del notebook EDA_Entregable1.ipynb: \dato{clave} muestra el valor
% definido en resultados_eda.tex (o una marca en rojo si no existe), y
% \sidato{clave}{texto}{alternativa} escribe el texto solo si la clave existe.
\makeatletter
\newcommand{\setdato}[2]{\@namedef{dato@#1}{#2}}
\newcommand{\dato}[2][]{\@ifundefined{dato@#2}{\pendiente{\ifx\relax#1\relax#2\else#1\fi}}{\@nameuse{dato@#2}}}
\newcommand{\sidato}[3]{\@ifundefined{dato@#1}{#3}{#2}}
\makeatother
\InputIfFileExists{resultados_eda.tex}{}{}
% \figura[ancho]{archivo}{texto}: incluye la figura del notebook si el PDF está
% en el proyecto; si no, muestra un recuadro en rojo con el texto.
\newcommand{\figura}[3][0.95]{\IfFileExists{figs/#2}{\includegraphics[width=#1\columnwidth]{figs/#2}}{\IfFileExists{#2}{\includegraphics[width=#1\columnwidth]{#2}}{\figpendiente{#3}}}}

\begin{document}

\title{Predicción de Crecimiento de Adopción en Protocolos Blockchain a partir de Fundamentales \emph{on-chain}}

\author{Michael~Stiven~Tabares~Tob\'on,~\IEEEmembership{Estudiante,~Universidad~de~Antioquia}\\
Mar\'ia~Fernanda~Atencia~Oliva,~\IEEEmembership{Estudiante,~Universidad~de~Antioquia}\\
Valeria~V\'asquez~Barrag\'an,~\IEEEmembership{Estudiante,~Universidad~de~Antioquia}\\
Jhon~Sebasti\'an~Mart\'inez~Soto,~\IEEEmembership{Estudiante,~Universidad~de~Antioquia}%
\thanks{Proyecto de aula para el curso Modelos y Simulación de Sistemas~II (Aprendizaje Supervisado), Departamento de Ingeniería de Sistemas, Universidad de Antioquia.}%
\thanks{Manuscrito recibido el 2 de octubre de 2026.}}

\markboth{Entregable I}{Descripción del problema}

\maketitle

\begin{abstract}
Este documento describe el problema de investigación del proyecto, consistente
en determinar si la información fundamental y el comportamiento histórico
reciente de un protocolo de finanzas descentralizadas (\emph{DeFi}), disponibles
hasta la semana $t$, contienen suficiente señal para anticipar si dicho protocolo
estará, durante la semana siguiente ($t+1$), entre el 10\,\% de protocolos con
mayor crecimiento relativo del valor total bloqueado (\emph{Total Value Locked},
TVL). Se describe la base de datos construida a partir de la API pública de
DeFiLlama, cuya unidad de observación es el par (\emph{protocolo, semana}):
\dato{N} observaciones de \dato{P} protocolos en \dato{T} semanas,
con una clase positiva de \dato{p}\,\%. Se detallan las variables, los datos
faltantes, la estrategia de imputación, la codificación y un análisis
exploratorio. El problema se formula como una clasificación binaria supervisada
con validación temporal, cuya métrica principal es el área bajo la curva
precisión--\emph{recall} (PR-AUC). En coherencia con la literatura especializada,
el diseño de características se orienta a los fundamentales de actividad
(\emph{Gross Merchandise Volume}, \emph{fees}) y a variables derivadas dinámicas
del TVL, y no solo a su nivel estático. Se presentan también las hipótesis del
estudio, las limitaciones de los datos y una revisión de cuatro trabajos de
aprendizaje de máquina sobre problemas afines.
\end{abstract}

\begin{IEEEkeywords}
Finanzas descentralizadas (DeFi), \emph{Total Value Locked} (TVL),
\emph{Gross Merchandise Volume} (GMV), clasificación supervisada, clases
desbalanceadas, fundamentales \emph{on-chain}, panel de datos, validación
temporal.
\end{IEEEkeywords}

% =====================================================================
%  I. DESCRIPCIÓN DEL PROBLEMA
% =====================================================================
\section{Descripción del problema}
\subsection{Contexto del problema y utilidad de una solución basada en \emph{machine learning}}

Las finanzas descentralizadas (\emph{DeFi}) constituyen un ecosistema de
servicios financieros ejecutados sobre contratos inteligentes, sin intermediarios
centralizados, en el que las métricas de actividad y de valor ---la más usada de
ellas el \emph{Total Value Locked} (TVL), que cuantifica los activos depositados
en los contratos de un protocolo--- son públicas y, en principio, verificables
\emph{on-chain}, aunque su cálculo en los agregadores depende de adaptadores
desarrollados por la comunidad y no está estandarizado \cite{saggese2025}. Este
ecosistema es altamente dinámico. El TVL agregado ha pasado por ciclos abruptos
de crecimiento y contracción (Fig.~\ref{fig:tvl}): cayó a
$\approx$\dato[40~000]{tvl2022}~M~USD a finales de 2022, subió a
$\approx$\dato[129~000]{tvl2025}~M~USD en 2025 y en 2026 llegó a estar
$\approx$\dato[39]{caida2026}\,\% por debajo de ese máximo. Además, el número
de protocolos elegibles del panel (Sec.~\ref{sec:datos}) pasó de
\dato[mín]{minE} a \dato[máx]{maxE} entre 2022 y 2025 (Fig.~\ref{fig:elegibles}).
En este entorno, el seguimiento manual de miles de protocolos es inviable, y
existe una asimetría de información considerable entre los desarrolladores de
los protocolos y sus usuarios \cite{ali2026}.

El problema consiste en determinar si la información fundamental y el
comportamiento histórico reciente de un protocolo ---disponibles hasta la semana
$t$--- contienen suficiente señal para anticipar si ese protocolo estará, durante
la semana siguiente ($t+1$), entre el 10\,\% de protocolos con mayor crecimiento
relativo del TVL de esa semana (definición formal en \eqref{eq:y}). La utilidad
de una solución basada en \emph{machine learning} (ML) es triple:

\begin{enumerate}
\item \textbf{Escalabilidad.} Un modelo supervisado permite procesar de forma
automática el panel completo de cada semana ---en la mediana,
\dato[med]{medE} protocolos elegibles (Fig.~\ref{fig:elegibles})---, algo
inabordable manualmente.
\item \textbf{Detección de relaciones no lineales.} La relación entre los
fundamentales \emph{on-chain} y el valor de un protocolo es no lineal y muy
volátil \cite{ali2026}, y en problemas afines los modelos de ML superan a las
referencias lineales y econométricas \cite{ghosh2023,jaquart2022}.
\item \textbf{Marco de evaluación riguroso.} La formulación como clasificación
binaria con validación temporal (Sec.~\ref{sec:paradigma}) permite responder de
forma verificable si existe señal predictiva en los fundamentales, y cuáles son
las variables más informativas.
\end{enumerate}

El proyecto es de naturaleza académica y experimental: no se busca construir un
sistema de \emph{trading} ni predecir precios de tokens, sino estudiar la
predictibilidad del crecimiento/adopción de protocolos a partir de fundamentales
\emph{on-chain}.

\subsection{Composición de la base de datos}
\label{sec:datos}

\textbf{Fuente de datos.} La base de datos se construye a partir de la API
pública y gratuita de DeFiLlama (\texttt{https://api.llama.fi}), cuyos endpoints
fueron verificados el 20/08/2026 \cite{defillama2026}:
{\raggedright
\begin{itemize}
\item \texttt{GET /protocols} --- catálogo de protocolos y categorías.
\item \texttt{GET /protocol/\{slug\}} --- TVL histórico diario (total y por
cadena).
\item \texttt{GET /summary/fees/\{protocol\}} --- \emph{fees}/revenue diarios.
\item \texttt{GET /summary/dexs/\{protocol\}} --- volumen DEX diario.
\item \texttt{GET /overview/fees} y \texttt{/overview/dexs} --- protocolos que
reportan \emph{fees} o volumen.
\item \texttt{GET /v2/historicalChainTvl} --- TVL agregado del ecosistema
(contexto de mercado).
\end{itemize}
\par}
Los datos se descargaron el \dato[dd/mm/2026]{fechadescarga}. DeFiLlama es también la
fuente de TVL de trabajos académicos recientes sobre DeFi
\cite{ali2026,metelski2022,grande2025}.

La granularidad nativa es diaria; el dataset se agrega a \textbf{semanas ISO}
(cierre en domingo UTC): el TVL semanal es el valor al cierre de la semana, y
los \emph{fees} y el volumen DEX se suman sobre los siete días. La unidad de
observación es el par \emph{(protocolo, semana)}: un protocolo aporta una fila
por cada semana en la que tiene datos válidos.

\textbf{Universo elegible y tamaño.} El universo elegible $\mathcal{E}_t$ de la
semana $t$ está restringido a protocolos con TVL semanal $\geq$ 1 M USD y con al
menos 8 semanas de historial previo; se excluyen los \emph{exchanges}
centralizados (categoría CEX), que DeFiLlama lista pero no son protocolos DeFi.
El período va del 3 de enero de 2022 a la última semana ISO completa disponible
en la descarga (\dato[dd/mm/2026]{ultimasemana}). El panel resultante tiene
\dato{N} observaciones de \dato{P} protocolos distintos en
\dato{T} semanas; el número de protocolos elegibles por semana varía entre
\dato[mín]{minE} y \dato[máx]{maxE}, con mediana \dato[med]{medE}
(Fig.~\ref{fig:elegibles}).

\textbf{Variable objetivo.} Para cada protocolo $i \in \mathcal{E}_t$ se define
el crecimiento logarítmico del TVL entre las semanas $t$ y $t+1$,
\begin{equation}
g_{i,t+1} = \ln\!\left(\frac{\text{TVL}_{i,t+1}}{\text{TVL}_{i,t}}\right),
\label{eq:g}
\end{equation}
y la etiqueta binaria
\begin{equation}
y_{i,t} = \begin{cases}
1, & \text{si } g_{i,t+1} \geq q_{0.9}(t+1),\\
0, & \text{en otro caso,}
\end{cases}
\label{eq:y}
\end{equation}
donde $q_{0.9}(t+1)$ es el percentil 90 de $\{g_{j,t+1} : j \in \mathcal{E}_t\}$.
La etiqueta es relativa a la cohorte de cada semana, de modo que no depende de si
todo el mercado sube o baja; por construcción, la clase positiva es cercana al
10\,\% (en la base final, \dato{p}\,\%). Si un protocolo deja de reportar
TVL en $t+1$, se le asigna $y_{i,t}=0$ y no entra en el cálculo del percentil,
para no descartar los casos de abandono.

\textbf{Fundamentales priorizados.} La literatura sobre valoración y predicción
de protocolos DeFi converge en señalar que las métricas de \textbf{actividad}
(distinguidas del capital inmóvil) son las más informativas. Metelski y Sobieraj
\cite{metelski2022} encuentran que el \emph{Gross Merchandise Volume} (GMV) ---el
volumen transaccional agregado de un protocolo, análogo al volumen DEX de
DeFiLlama--- es la \emph{única} variable que causa, en el sentido de Granger, la
valoración futura, mientras que el TVL no la anticipa: es la valoración la que
anticipa al TVL. En el mismo sentido, Ali \cite{ali2026}, mediante análisis SHAP,
encuentra que el GMV ($\approx 0.062$) supera al TVL ($\approx 0.050$) como
determinante de la valoración; le siguen la clasificación del protocolo como DEX
($\approx 0.039$), la inflación de tokens ($\approx 0.034$, con efecto negativo
por dilución) y los ingresos totales ($\approx 0.033$). Fuera de DeFi, Gu
\emph{et al.} \cite{gukellyxiu2020} concluyen que los predictores dominantes del
retorno de sección cruzada son el \emph{momentum}, la liquidez y la volatilidad:
variables dinámicas, no niveles estáticos. Un informe de industria de Greenfield
Capital \cite{greenfield2025} matiza este resultado: con \emph{Random Forest}
sobre 77 protocolos, encuentra que desde 2024 el TVL y los \emph{fees} tienen
importancias similares en el horizonte de 6 meses ($\approx$19.7\,\% y
$\approx$19.2\,\%), por encima del volumen DEX ($\approx$11.8\,\%), aunque con
un diseño contemporáneo, no predictivo.

Por ello, el panel incorpora (i) características de actividad ---el GMV/volumen
DEX y \emph{fees} semanales, su crecimiento a 1 y 4 semanas y su razón respecto
al TVL (eficiencia de monetización y utilización del capital)--- y (ii) variables
derivadas del TVL que añaden dinamismo: crecimiento a 1/2/4/8/12/26 semanas,
\emph{momentum} a 2/4/12/26 semanas, aceleración y volatilidad a 4/8/12/26
semanas, junto con la posición relativa (\emph{rank}) del TVL dentro de la semana.
En total son 24 variables numéricas (Tabla~\ref{tab:vars}), a las que se suman
2 variables indicadoras de disponibilidad y la categoría del protocolo. Todas se
calculan con información disponible hasta el cierre de la semana $t$.

% --- Tabla de variables predictoras ------------------------------------
\begin{table}[!t]
\caption{Variables predictoras: actividad (GMV, \emph{fees}, razones) y
variables derivadas dinámicas del TVL. Todas se calculan con información
$\leq$ cierre de la semana $t$.\label{tab:vars}}
\centering
\footnotesize
\setlength{\tabcolsep}{3pt}
\begin{tabular}{@{}L{2.75cm}cL{4.8cm}@{}}
\toprule
\textbf{Variable} & \textbf{Cant.} & \textbf{Definición (tipo)} \\
\midrule
\texttt{log\_fees\_week} & 1 & $\ln(1+\textit{fees}$ semanales$)$; actividad (continua) \\
\texttt{g\_fees\_1w}, \texttt{g\_fees\_4w} & 2 & crecimiento de \emph{fees} a 1 y 4 semanas (continua) \\
\texttt{log\_dex\_vol\_week} & 1 & $\ln(1+$volumen DEX$)$ $\approx$ GMV; actividad (continua) \\
\texttt{ratio\_fees\_tvl} & 1 & \emph{fees}/TVL: eficiencia de monetización (continua) \\
\texttt{ratio\_vol\_tvl} & 1 & volumen DEX/TVL: utilización del capital (continua) \\
\texttt{fees\_tvl\_rank\_w} & 1 & percentil de \emph{fees}/TVL dentro de la semana (percentil) \\
\addlinespace
\texttt{g\_tvl\_1w...26w} & 6 & crecimiento del TVL a 1, 2, 4, 8, 12 y 26 semanas (continua) \\
\texttt{mom\_2w...26w} & 4 & \emph{momentum}: media del crecimiento semanal en 2, 4, 12 y 26 semanas (continua) \\
\texttt{accel}, \texttt{accel\_12w} & 2 & aceleración del crecimiento (continua) \\
\texttt{vol\_4w...26w} & 4 & volatilidad: desv. estándar del crecimiento en 4, 8, 12 y 26 semanas (continua) \\
\texttt{tvl\_rank\_w} & 1 & percentil del TVL dentro de la semana (percentil) \\
\midrule
\multicolumn{3}{@{}L{8.4cm}@{}}{Total: 24 variables numéricas. Además: 2 indicadoras (\texttt{has\_fees}, \texttt{has\_dex\_vol}) y la categoría del protocolo en \emph{one-hot}\sidato{ncat}{ (\dato{ncat} columnas)}{}.} \\
\bottomrule
\end{tabular}
\end{table}

\textbf{Datos faltantes e imputación.} No todos los protocolos reportan
\emph{fees}, y el volumen DEX solo existe para los protocolos de intercambio; por
eso estas variables concentran los faltantes (Tabla~\ref{tab:eda}). La estrategia de imputación es:
(i) los \emph{fees} o el volumen no reportados se imputan como 0 y se agregan las
indicadoras \texttt{has\_fees} y \texttt{has\_dex\_vol}, para distinguir ``no
reporta'' de ``reporta cero''; (ii) los huecos puntuales en la serie diaria de
TVL se rellenan hacia adelante hasta 7 días, sin rellenar nunca hacia atrás
---a diferencia de \cite{ali2026}---, para no introducir información futura;
(iii) cuando una ventana es más larga que el historial disponible (por ejemplo,
26 semanas en un protocolo con 10), se imputa la mediana transversal de la
semana, como en \cite{gukellyxiu2020}; y (iv) los protocolos con menos de 8
semanas de historial quedan fuera por el criterio de elegibilidad.

\textbf{Codificación.} Los montos (\emph{fees}, volumen) se transforman con
$\ln(1+x)$ por su fuerte asimetría, y los crecimientos se calculan como
diferencias de $\ln(1+x)$, lo que los deja definidos aun cuando el valor previo
es 0. Las posiciones relativas se codifican como percentiles dentro de la
semana, en $[0,1]$. Las variables continuas se recortan en los percentiles 1 y
99 de cada semana para limitar valores extremos, sin modificar la variable
objetivo. La categoría del protocolo (DEX, \emph{Lending}, \emph{Liquid Staking},
etc.) se codifica en \emph{one-hot}, agrupando en ``Otros'' las categorías con
menos de \dato{m} protocolos\sidato{ncat}{, lo que da \dato{ncat} columnas
binarias; en total, cada observación se describe con
\the\numexpr 26+\csname dato@ncat\endcsname\relax{} columnas}{}. La etiqueta es
binaria (0/1), y el escalamiento estándar que requieren modelos como SVM o redes
neuronales se ajustará solo con el conjunto de entrenamiento.

\textbf{Análisis exploratorio.} La Tabla~\ref{tab:eda} resume las variables
principales y su porcentaje de faltantes; la Fig.~\ref{fig:eda1} muestra la
evolución del TVL agregado y del número de protocolos elegibles, y la
Fig.~\ref{fig:eda2}, la distribución del crecimiento semanal y la correlación
entre predictores. \dato[Describir en 2--3 frases los hallazgos: asimetría
de las variables de actividad, magnitud del umbral $q_{0.9}$, estabilidad del
balance de clases por semana y pares de variables muy correlacionadas.]{hallazgos}
Un umbral $q_{0.9}$ negativo corresponde a semanas de caída generalizada, en las
que incluso el decil superior pierde TVL; aun así, la etiqueta relativa conserva
cerca de 10\,\% de positivos en cualquier régimen de mercado. Además,
\texttt{fees\_tvl\_rank\_w} es el percentil semanal de
\texttt{ratio\_fees\_tvl}, por lo que ese par es casi redundante por
construcción.

% --- Tabla de estadísticas descriptivas --------------------------------
\begin{table}[!t]
\caption{Estadísticas descriptivas y porcentaje de datos faltantes (antes de
imputar) de variables seleccionadas\label{tab:eda}}
\centering
\footnotesize
\setlength{\tabcolsep}{3.5pt}
\begin{tabular}{@{}lcccc@{}}
\toprule
\textbf{Variable} & \textbf{Media} & \textbf{Mediana} & \textbf{Desv.} & \textbf{\% falt.} \\
\midrule
$g_{i,t+1}$ (objetivo) & \dato[--]{t-g-media} & \dato[--]{t-g-mediana} & \dato[--]{t-g-desv} & \dato[--]{t-g-falt} \\
\texttt{g\_tvl\_1w} & \dato[--]{t-gtvl1-media} & \dato[--]{t-gtvl1-mediana} & \dato[--]{t-gtvl1-desv} & \dato[--]{t-gtvl1-falt} \\
\texttt{log\_fees\_week} & \dato[--]{t-logfees-media} & \dato[--]{t-logfees-mediana} & \dato[--]{t-logfees-desv} & \dato[--]{t-logfees-falt} \\
\texttt{log\_dex\_vol\_week} & \dato[--]{t-logdex-media} & \dato[--]{t-logdex-mediana} & \dato[--]{t-logdex-desv} & \dato[--]{t-logdex-falt} \\
\texttt{ratio\_fees\_tvl} & \dato[--]{t-ratiofees-media} & \dato[--]{t-ratiofees-mediana} & \dato[--]{t-ratiofees-desv} & \dato[--]{t-ratiofees-falt} \\
\texttt{vol\_12w} & \dato[--]{t-vol12-media} & \dato[--]{t-vol12-mediana} & \dato[--]{t-vol12-desv} & \dato[--]{t-vol12-falt} \\
\midrule
\multicolumn{5}{@{}l}{Clase positiva ($y_{i,t}=1$): \dato{p}\,\% de las observaciones.} \\
\bottomrule
\end{tabular}
\end{table}

% --- Figuras del análisis exploratorio ---------------------------------
% Las figuras las genera el notebook EDA_Entregable1.ipynb; \figura las incluye
% cuando el PDF está en el proyecto (en la raíz o en la carpeta figs/).
\begin{figure}[!t]
\centering
\subfloat[TVL agregado del ecosistema por semana (eje $y$: miles de millones de USD).\label{fig:tvl}]{%
\figura{tvl_agregado.pdf}{Eje $x$: semana (2022--2026); eje $y$: TVL agregado (miles de millones de USD). Fuente: \texttt{/v2/historicalChainTvl}.}}
\\
\subfloat[Protocolos del universo elegible $\mathcal{E}_t$ por semana (eje $y$: número de protocolos).\label{fig:elegibles}]{%
\figura{protocolos_elegibles.pdf}{Eje $x$: semana; eje $y$: número de protocolos elegibles.}}
\caption{Composición temporal de la base de datos.}
\label{fig:eda1}
\end{figure}

\begin{figure}[!t]
\centering
\subfloat[Distribución de $g_{i,t+1}$ con el umbral promedio $q_{0.9}$; se omite el 1\,\% de valores más extremos (eje $y$: frecuencia, en escala logarítmica).\label{fig:dist}]{%
\figura{distribucion_crecimiento.pdf}{Histograma. Eje $x$: crecimiento logarítmico semanal del TVL; eje $y$: frecuencia; línea vertical en $q_{0.9}$.}}
\\
\subfloat[Correlación de Spearman entre las 24 variables numéricas.\label{fig:corr}]{%
\figura[1]{correlacion.pdf}{Mapa de calor. Ejes $x$ e $y$: variables de la Tabla~\ref{tab:vars}; color: coeficiente de Spearman.}}
\caption{Variable objetivo y relación entre predictores.}
\label{fig:eda2}
\end{figure}

\textbf{Limitaciones de los datos.} El TVL se calcula como la cantidad de tokens
depositados multiplicada por su precio, de modo que parte de su crecimiento
refleja variaciones de precio y no nuevos depósitos; la etiqueta relativa por
semana atenúa el efecto general del mercado, pero no la distinta exposición de
cada protocolo a activos como ETH. Además, el cálculo depende de adaptadores
comunitarios sin estandarizar, algunos de los cuales consultan servidores
externos, y un mismo depósito puede contarse en más de un protocolo
\cite{saggese2025}. Por último, los protocolos que dejan de operar dejan de
reportar datos, lo que puede introducir sesgo de supervivencia.

\subsection{Paradigma de aprendizaje}
\label{sec:paradigma}

El problema se aborda con aprendizaje \textbf{supervisado}, como una tarea de
\textbf{clasificación binaria} sobre el panel (protocolo, semana): a partir de
las variables de la semana $t$ (Tabla~\ref{tab:vars}) se predice la etiqueta
$y_{i,t}$ de \eqref{eq:y}. Se prefiere la clasificación a la regresión del
crecimiento porque $g_{i,t+1}$ tiene colas muy pesadas (Fig.~\ref{fig:dist}), y
predecir si un protocolo estará en el decil superior es más estable y responde
directamente la pregunta de investigación. Gu \emph{et al.}
\cite{gukellyxiu2020} usan una lógica similar al ordenar los activos en deciles
según su predicción, y Jaquart \emph{et al.} \cite{jaquart2022} formulan una
clasificación binaria análoga: si cada criptomoneda superará la mediana de
rendimiento del día siguiente.

Como la clase positiva representa cerca del 10\,\% de las observaciones, la
métrica principal es el área bajo la curva precisión--\emph{recall} (PR-AUC),
estimada como la precisión promedio
\begin{equation}
\text{AP} = \sum_{n} (R_n - R_{n-1})\, P_n,
\label{eq:ap}
\end{equation}
donde $P_n$ y $R_n$ son la precisión y el \emph{recall} en el $n$-ésimo umbral.
A diferencia del AUC-ROC, la PR-AUC no se ve dominada por la clase mayoritaria, y
un clasificador aleatorio obtiene un valor igual a la prevalencia ($\approx 0.10$).
El desbalance se tratará en el Entregable~II con ponderación de clases o técnicas
de sub y sobremuestreo, aplicadas solo al conjunto de entrenamiento.

La validación es temporal, sin mezclar semanas futuras en el entrenamiento:
entrenamiento de 2022 a 2024, validación en 2025 y prueba en 2026\sidato{nentr}{
(\dato{nentr}, \dato{nval} y \dato{nprueba} observaciones, respectivamente)}{}.
Como la etiqueta de la semana $t$ usa el TVL de $t+1$, se descarta la última semana de
cada bloque para evitar solapamiento entre conjuntos. El año de prueba coincide
con una fase de contracción del TVL (Fig.~\ref{fig:tvl}), lo que permitirá
evaluar la robustez de los modelos ante un cambio de régimen.

\subsection{Hipótesis}
\label{sec:hipotesis}

\begin{itemize}
\item \textbf{H1:} las variables de actividad (\emph{fees}, volumen DEX y sus
razones con el TVL) aportan más poder predictivo que las variables basadas solo
en el nivel del TVL.
\item \textbf{H2:} los modelos no lineales (ensambles de árboles, redes
neuronales) superan en PR-AUC a un modelo lineal de referencia (regresión
logística).
\end{itemize}
H1 se evaluará comparando modelos entrenados con cada grupo de variables, y H2,
comparando los modelos del Entregable~II contra la regresión logística.

% =====================================================================
%  II. ESTADO DEL ARTE  (la guía pide no superar una página)
% =====================================================================
\section{Estado del arte}
\label{sec:estado}

Se revisaron artículos de revista (2020--2026) de \emph{machine learning} que
usaran la misma fuente de datos (DeFiLlama) o variables equivalentes (TVL, GMV,
\emph{fees}), o que resolvieran el mismo tipo de problema: predecir el desempeño
relativo de muchos activos. No encontramos trabajos de ML que usen DeFiLlama
para predecir el crecimiento del TVL, que es el vacío que aborda este proyecto.

\subsection{Trabajos revisados}

En la Tabla~\ref{tab:sota}, Ali \cite{ali2026} es el trabajo más cercano en
datos (DeFiLlama y las mismas familias de variables), aunque predice la
valoración y no el crecimiento del TVL; Ghosh \emph{et al.} \cite{ghosh2023}
muestran que los ensambles de árboles predicen bien los precios de tokens DeFi;
y Gu \emph{et al.} \cite{gukellyxiu2020} y Jaquart \emph{et al.}
\cite{jaquart2022} son los más cercanos en estructura: predicen el desempeño
relativo en un panel de muchos activos, y los segundos lo formulan, como
nosotros, como una clasificación binaria respecto a la cohorte de cada periodo.

% --- Tabla del estado del arte (ancho de página) -----------------------
\begin{table*}[!t]
\caption{Trabajos de aprendizaje de máquina revisados\label{tab:sota}}
\centering
\footnotesize
\setlength{\tabcolsep}{3pt}
\renewcommand{\arraystretch}{1.05}
\begin{tabular}{@{}L{0.075\textwidth}L{0.095\textwidth}L{0.29\textwidth}L{0.17\textwidth}L{0.085\textwidth}L{0.21\textwidth}@{}}
\toprule
\textbf{Trabajo} & \textbf{Paradigma} & \textbf{Técnica y datos} & \textbf{Validación} & \textbf{Métricas} & \textbf{Resultados} \\
\midrule
Ali (2026) \cite{ali2026} & Regresión supervisada (valoración) &
\emph{Super Learner} apilado (\emph{Extra Trees}, CatBoost y SVR; KNN como meta-aprendiz) con Optuna. 30 protocolos DeFi, 5~370 observaciones diarias de 2022 (DeFiLlama, TokenTerminal y DappRadar). &
Validación cruzada de 10 pliegues para elegir algoritmos; validación cruzada temporal de panel con ventana expandida. &
RMSE, MAE, $R^2$; Wilcoxon. &
RMSE $=0.085$, MAE $=0.065$, $R^2=0.97$; error 25--36\,\% menor que el de los modelos alternativos ($p<0.01$). SHAP: GMV $>$ TVL $>$ clase DEX. \\
\addlinespace
Ghosh \emph{et al.} (2023) \cite{ghosh2023} & Regresión supervisada (precio diario) &
ISOMAP + \emph{Gradient Boosting} y UMAP + \emph{Random Forest} con indicadores técnicos, macroeconómicos y de medios. Precios diarios de 4 tokens DeFi y 4 NFT, 2020--2022. &
Partición 85/15; \emph{GridSearchCV}; comparación con DT, SVR, ARIMA y SARIMA (Kruskal-Wallis). &
NSE, TI, IA, DA. &
En prueba, NSE de 0.85 a 0.98 y DA de 0.89 a 0.96 en los tokens DeFi; superan a los modelos de comparación. \\
\addlinespace
Gu \emph{et al.} (2020) \cite{gukellyxiu2020} & Regresión supervisada (retorno mensual) &
OLS, ENet, PCR, PLS, GLM, RF, GBRT y redes neuronales de 1 a 5 capas; cerca de 30~000 acciones y 920 predictores, 1957--2016. &
Entrenamiento 1957--1974, validación 1975--1986 y prueba 1987--2016; reentrenamiento anual. &
$R^2_{oos}$; Diebold-Mariano; Sharpe. &
Árboles y redes: $R^2_{oos}$ mensual de 0.33 a 0.40\,\% (OLS con 3 variables: 0.16\,\%); cartera larga-corta por deciles con Sharpe de 1.35. \\
\addlinespace
Jaquart \emph{et al.} (2022) \cite{jaquart2022} & Clasificación binaria supervisada (desempeño relativo) &
LSTM, GRU, TCN, GBC, RF y regresión logística sobre los retornos de los 90 días previos; 100 criptomonedas de mayor capitalización (CoinGecko), 2018--2022. Etiqueta: retorno del día siguiente $\geq$ mediana. &
5 periodos móviles de 800 días: entrenamiento (500), validación de hiperparámetros (150) y prueba (150). &
Exactitud; prueba binomial; Sharpe. &
Exactitud de 52.9\,\% (LR) a 54.2\,\% (RF); 57.5--59.5\,\% en el 10\,\% de predicciones más confiables. Sharpe tras costos: 3.23 (LSTM) frente a 1.33 del mercado. \\
\bottomrule
\end{tabular}
\end{table*}

\subsection{Definición de métricas}
\label{sec:met}

RMSE, MAE, $R^2$, exactitud, precisión, \emph{recall}, F1 y AUC-ROC son
métricas estándar que no se redefinen aquí, y la \mbox{PR-AUC} se definió en
\eqref{eq:ap}. Para las demás, $y_t$ es el
valor real, $\hat{y}_t$ el predicho, $\bar{y}$ la media real y $N$ el número de
observaciones. Gu \emph{et al.} \cite{gukellyxiu2020} usan un $R^2$ fuera de
muestra que compara contra una predicción de cero, no contra la media histórica:
\begin{equation}
R^2_{oos} = 1 - \frac{\sum (r_{i,t+1} - \hat{r}_{i,t+1})^2}{\sum r_{i,t+1}^2}.
\label{eq:r2oos}
\end{equation}
Ghosh \emph{et al.} \cite{ghosh2023} usan la eficiencia de Nash-Sutcliffe (NSE),
el índice de Theil (TI) y el índice de concordancia (IA):
\begin{align}
\text{NSE} &= 1 - \frac{\sum (\hat{y}_t - y_t)^2}{\sum (y_t - \bar{y})^2}, \label{eq:nse}\\
\text{TI} &= \frac{\sqrt{\tfrac{1}{N}\sum (\hat{y}_t - y_t)^2}}{\sqrt{\tfrac{1}{N}\sum \hat{y}_t^2} + \sqrt{\tfrac{1}{N}\sum y_t^2}}, \label{eq:ti}\\
\text{IA} &= 1 - \frac{\sum (\hat{y}_t - y_t)^2}{\sum \left(|\hat{y}_t - \bar{y}| + |y_t - \bar{y}|\right)^2}, \label{eq:ia}
\end{align}
y la exactitud direccional $\text{DA} = \frac{1}{N}\sum_t D_t$, con $D_t = 1$ si
$(y_{t+1} - y_t)(\hat{y}_{t+1} - y_t) \geq 0$ y $D_t = 0$ en otro caso. NSE e IA
cercanos a 1 y TI cercano a 0 indican mejores predicciones. La razón de Sharpe
\cite{gukellyxiu2020,jaquart2022} es $\text{SR} = \bar{r}_p/\sigma_p$: la media
del retorno (en exceso) de una cartera dividida por su desviación estándar. Las
pruebas de Wilcoxon, Kruskal-Wallis y Diebold-Mariano contrastan si las
diferencias de error entre modelos son significativas, y la binomial, si la
exactitud supera el 50\,\% del azar.

\subsection{Síntesis y contribución al problema}

La literatura apunta en tres direcciones: la señal útil está en la dinámica y la
actividad (\emph{momentum}, volatilidad, \emph{fees}, volumen) más que en el
nivel del TVL; los métodos no lineales superan a los lineales, aunque la señal es
débil ---en \cite{jaquart2022}, la exactitud pasa de 52.9\,\% con regresión
logística a 54.2\,\% con RF, frente al 50\,\% del azar---, por lo que la \mbox{PR-AUC}
se comparará con la de un clasificador aleatorio y la de la regresión logística;
y la validación debe respetar el orden temporal. Tres de los cuatro trabajos
formulan una regresión; nuestro proyecto sigue la lógica de \cite{jaquart2022}
con una clasificación relativa: ¿estará el protocolo en el decil superior de
crecimiento la próxima semana?

% =====================================================================
%  Bibliografía (BibTeX, estilo IEEE)
% =====================================================================
\bibliographystyle{IEEEtran}
\bibliography{references}

\end{document}
